\documentclass[10pt,a4paper]{article}
\usepackage[utf8]{inputenc}
\usepackage[T1]{fontenc}
\usepackage{amsmath,amssymb}
\usepackage[margin=1.8cm]{geometry}
\usepackage[table]{xcolor}
\usepackage{graphicx}
\usepackage{caption}

\newcommand{\dt}{\Delta t}
\captionsetup{font=small}
\setlength{\parskip}{0.3em}
\setlength{\abovedisplayskip}{5pt}
\setlength{\belowdisplayskip}{5pt}
\setlength{\abovedisplayshortskip}{3pt}
\setlength{\belowdisplayshortskip}{3pt}
\pagestyle{empty}

\begin{document}

\begin{center}
{\large\bfseries The conductance step: closed-form integration}\\[0.2em]
{\small branch \texttt{feat/flyvis-flow-conductance-E} \quad|\quad 16 September 2026}
\end{center}

\textbf{The model.} For postsynaptic neuron $i$, with $g_{ij}\ge 0$ the
conductance of edge $j\!\to\!i$, $E_{ij}$ the reversal that edge drives toward,
$E_i$ the leak reversal of the cell itself (flyvis stores it as the parameter
\texttt{bias}, one learned value per cell type), $\mathrm{stim}_i$ the external
stimulus current and $\tau_i$ the membrane time constant,
\begin{equation}
\tau_i\frac{dv_i}{dt}
= -v_i + E_i + \sum_j g_{ij}\,u_j\,(E_{ij}-v_i) + \mathrm{stim}_i ,
\qquad u_j \equiv \mathrm{relu}\bigl(v_j(t)\bigr).
\label{eq:model}
\end{equation}
$u_j$ is the presynaptic drive. It is a \emph{known number} at the moment the step
is taken --- the state of the other neurons at time $t$ --- not an unknown of the
equation for $v_i$, which is what makes the closed form below available.

\textbf{It is affine in $v_i$.} Define the conductance-weighted sum of the
reversals reaching $i$, and the total synaptic conductance onto $i$ in units of
its own leak conductance:
\begin{equation}
S_i=\sum_j g_{ij}\,u_j\,E_{ij},
\qquad
G_i=\sum_j g_{ij}\,u_j\ \ (\ge 0).
\end{equation}
Both are known at time $t$. The synaptic sum then splits into a part independent
of $v_i$ and a part proportional to it,
\begin{equation}
\sum_j g_{ij}\,u_j\,(E_{ij}-v_i)
=\underbrace{\sum_j g_{ij}u_jE_{ij}}_{S_i}
-\;v_i\underbrace{\sum_j g_{ij}u_j}_{G_i}
= S_i - v_i G_i ,
\end{equation}
so the right-hand side of \eqref{eq:model} collects the two $v_i$ terms --- the
leak $-v_i$ and the shunt $-v_iG_i$ --- into one:
\begin{align}
\tau_i\frac{dv_i}{dt}
&= -v_i + E_i + S_i - v_iG_i + \mathrm{stim}_i \nonumber\\[0.15em]
&= -(1+G_i)\,v_i + \bigl(E_i+S_i+\mathrm{stim}_i\bigr) \nonumber\\[0.15em]
&= (1+G_i)\left[\frac{E_i+S_i+\mathrm{stim}_i}{1+G_i}-v_i\right] .
\end{align}
The bracket is the distance to the voltage at which the right-hand side vanishes,
so \eqref{eq:model} is a pure relaxation,
\begin{equation}
\tau_i\frac{dv_i}{dt}=(1+G_i)\bigl(v^{\infty}_i-v_i\bigr),
\qquad
v^{\infty}_i=\frac{E_i+S_i+\mathrm{stim}_i}{1+G_i},
\qquad
\tau^{\mathrm{eff}}_i=\frac{\tau_i}{1+G_i}.
\label{eq:relax}
\end{equation}

\textbf{Integration over one step.} $G_i$, $S_i$ and $\mathrm{stim}_i$ are read off the state
at time $t$ and held there across the step, as forward Euler also does. Then $v^{\infty}_i$ and $\tau^{\mathrm{eff}}_i$ are constants over
$[t,t+\dt]$, \eqref{eq:relax} separates, and the integral is elementary:
\begin{equation}
\int_{v_i(t)}^{v_i(t+\dt)}\frac{dv_i}{v^{\infty}_i-v_i}
=\int_{t}^{t+\dt}\frac{dt'}{\tau^{\mathrm{eff}}_i}
\qquad\Longrightarrow\qquad
\ln\frac{v^{\infty}_i-v_i(t)}{v^{\infty}_i-v_i(t+\dt)}
=\frac{\dt}{\tau^{\mathrm{eff}}_i}\;\equiv\;z ,
\end{equation}
so that, with $z=(1+G_i)\dt/\tau_i$ the number of effective relaxation times one
step covers,
\begin{equation}
\boxed{\;v_i(t+\dt)=v^{\infty}_i+\bigl(v_i(t)-v^{\infty}_i\bigr)\,e^{-z}\;}
\label{eq:exact}
\end{equation}

\end{document}
